\section{Copy constraints in PLONK}
\label{copy-constraints}

Now we elaborate on Step 3 which we deferred back in
Section \ref{copy-constraint-deferred}. As an example, the constraints might be:
\[a_{1} = a_{4} = c_{4}\quad\text{and }\quad b_{2} = c_{1.}\] Before we
show how to check this, we provide a solution to a ``simpler'' problem
called ``permutation check''. Then we explain how to deal with the full
copy check.

\subsection{Easier case: permutation check}

\begin{problem}{}{}
Suppose we have polynomials $P,Q \in \FF_{q}[X]$
which encode two vectors of values \[\begin{aligned}
\overset{\rightarrow}{p} & = \langle P( \omega^{1} ),P( \omega^{2} ),\ldots,P( \omega^{n} )\rangle \\
\overset{\rightarrow}{q} & = \langle Q( \omega^{1} ),Q( \omega^{2} ),\ldots,Q( \omega^{n} )\rangle.
\end{aligned}\] Is there a way that one can quickly verify
$\overset{\rightarrow}{p}$ and $\overset{\rightarrow}{q}$ are the
same up to permutation of the $n$ entries?
\end{problem}


Well, actually, it would be necessary and sufficient for the identity
\begin{eqnarray}
\label{permcheck-poly}
( T + P( \omega^{1} ) )( T + P( \omega^{2} ) )\cdots( T + P( \omega^{n} ) ) \nonumber \\
 = ( T + Q( \omega^{1} ) )( T + Q( \omega^{2} ) )\cdots( T + Q( \omega^{n} ) )
\end{eqnarray}
to be true, in the sense that both sides are the same polynomial in a
single formal variable $T$. And for that, it is sufficient that a
single random challenge $T = \lambda$ passes
Equation \eqref{permcheck-poly}: If the two sides of
Equation \eqref{permcheck-poly} are not the same
polynomial, then the two sides can have at most $n - 1$ common values.
So for a randomly chosen $\lambda$ (chosen from a field with
$q \approx 2^{256}$ elements), the chances that $T = \lambda$ passes
Equation \eqref{permcheck-poly} are extremely small.

We can then get a proof of
Equation \eqref{permcheck-poly} using the technique of
adding an \emph{accumulator polynomial}. The idea is this: Victor picks
a random challenge $\lambda \in \FF_{q}$. Peggy then
interpolates the polynomial
$F_{P} \in \FF_{q}[T]$ such that
\[\begin{aligned}
F_{P}( \omega^{1} ) & = \lambda + P( \omega^{1} ) \\
F_{P}( \omega^{2} ) & = ( \lambda + P( \omega^{1} ) )( \lambda + P( \omega^{2} ) ) \\
 & \vdots \\
F_{P}( \omega^{n} ) & = ( \lambda + P( \omega^{1} ) )( \lambda + P( \omega^{2} ) )\cdots( \lambda + P( \omega^{n} ) ).
\end{aligned}\] Then the accumulator
$F_{Q} \in \FF_{q}[T]$ is defined analogously.

So to prove Equation \eqref{permcheck-poly}, the
following algorithm works:

\begin{algo}{Permutation check}{}
Suppose Peggy has committed \(\operatorname{Com}(P)\)
and \(\operatorname{Com}(Q)\).

\begin{enumerate}
\item
  Victor sends a random challenge $\lambda \in \FF_{q}$.
\item
  Peggy interpolates polynomials $F_{P}[T]$ and
  $F_{Q}[T]$ such that
  \(F_{P}( \omega^{k} ) = \prod_{i \leq k}( \lambda + P( \omega^{i} ) )\).
  Define $F_{Q}$ similarly. Peggy sends \(\operatorname{Com}(F_{P})\)
  and \(\operatorname{Com}(F_{Q})\).
\item
  Peggy uses Algorithm \ref{tcb:root-check} to prove all of the following statements:

  \begin{itemize}
  \item
    \(F_{P}(X) - ( \lambda + P(X) )\) vanishes at
    $X = \omega$;
  \item
    \(F_{P}(\omega X) - ( \lambda + P(\omega X) )F_{P}(X)\)
    vanishes at $X \in \{ \omega,\ldots,\omega^{n - 1} \}$;
  \item
    The previous two statements also hold with $F_{P}$ replaced by
    $F_{Q}$;
  \item
    \(F_{P}(X) - F_{Q}(X)\) vanishes at $X = 1$.
  \end{itemize}
\end{enumerate}
\end{algo}

\subsection{Copy check}

Moving on to copy check, let us look at a concrete example where
$n = 4$. Suppose that our copy constraints were
\[ \redoval{(a_{1})} =
\redoval{(a_{4})} =
\redoval{(c_{4})}
\quad\text{and }\quad \bluerect{(b_{2})} = \bluerect{(c_{1})}.\]
(We have colored the variables that will move around for
readability.) So, the copy constraint means we want the following
equality of matrices:
\begin{equation}
\label{copy1}
\begin{pmatrix}
a_{1} & b_{1} & c_{1} \\
a_{2} & b_{2} & c_{2} \\
a_{3} & b_{3} & c_{3} \\
a_{4} & b_{4} & c_{4}
\end{pmatrix} = \begin{pmatrix}
\redoval{(a_{4})} & b_{1} & \bluerect{(b_{2})} \\
a_{2} & \bluerect{(c_{1})} & c_{2} \\
a_{3} & b_{3} & c_{3} \\
\redoval{(c_{4})} & b_{4} & \redoval{(a_{1})}
\end{pmatrix}.
\end{equation}

Again, our goal is to make this into a \emph{single} equation. There is
a really clever way to do this by tagging each entry with
$+ \eta^{j}\omega^{k}\mu$ in reading order for $j = 0,1,2$ and
$k = 1,\ldots,n$; here $\eta \in \FF_{q}$ is any number
such that $\eta^{2}$ does not happen to be a power of $\omega$, so
all the tags are distinct. Specifically, if
Equation \eqref{copy1} is true, then for any
$\mu \in \FF_{q}$, we also have
\begin{equation}
\label{copy2}
\begin{aligned}
 & \begin{pmatrix}
a_{1} + \omega^{1}\mu & b_{1} + \eta\omega^{1}\mu & c_{1} + \eta^{2}\omega^{1}\mu \\
a_{2} + \omega^{2}\mu & b_{2} + \eta\omega^{2}\mu & c_{2} + \eta^{2}\omega^{2}\mu \\
a_{3} + \omega^{3}\mu & b_{3} + \eta\omega^{3}\mu & c_{3} + \eta^{2}\omega^{3}\mu \\
a_{4} + \omega^{4}\mu & b_{4} + \eta\omega^{4}\mu & c_{4} + \eta^{2}\omega^{4}\mu
\end{pmatrix}  \\
 = & \begin{pmatrix}
\redoval{(a_{4})} + \omega^{1}\mu & b_{1} + \eta\omega^{1}\mu & \bluerect{(b_{2})} + \eta^{2}\omega^{1}\mu \\
a_{2} + \omega^{2}\mu & \bluerect{(c_{1})} + \eta\omega^{2}\mu & c_{2} + \eta^{2}\omega^{2}\mu \\
a_{3} + \omega^{3}\mu & b_{3} + \eta\omega^{3}\mu & c_{3} + \eta^{2}\omega^{3}\mu \\
\redoval{(c_{4})} + \omega^{4}\mu & b_{4} + \eta\omega^{4}\mu & \redoval{(a_{1})} + \eta^{2}\omega^{4}\mu
\end{pmatrix}.
\end{aligned}
\end{equation}

Now how can the prover establish Equation \eqref{copy2}
succinctly? The answer is to run a permutation check on the $3n$
entries of Equation \eqref{copy2}! The prover will simply prove
that the twelve matrix entries of the matrix on the left are a
permutation of the twelve matrix entries of the matrix on the right.

The reader should check that this is correct! If the prover starts with
values $a_{i}$, $b_{i}$, and $c_{i}$ that do not satisfy all the
copy constraints, then a randomly selected $\mu$ is very unlikely to
satisfy this permutation check. The right-hand side will not be a
permutation of the left-hand side, and the check will fail.

To clean things up, shuffle the $12$ terms on the right-hand side of
Equation \eqref{copy2} so that each variable is in the cell it
started at. We want to prove
\begin{equation}
\label{copy3}
\begin{array}{c}
\begin{pmatrix}
a_{1} + \omega^{1}\mu & b_{1} + \eta\omega^{1}\mu & c_{1} + \eta^{2}\omega^{1}\mu \\
a_{2} + \omega^{2}\mu & b_{2} + \eta\omega^{2}\mu & c_{2} + \eta^{2}\omega^{2}\mu \\
a_{3} + \omega^{3}\mu & b_{3} + \eta\omega^{3}\mu & c_{3} + \eta^{2}\omega^{3}\mu \\
a_{4} + \omega^{4}\mu & b_{4} + \eta\omega^{4}\mu & c_{4} + \eta^{2}\omega^{4}\mu
\end{pmatrix} \\ \\
\text{is a permutation of } \\ \\
\begin{pmatrix}
a_{1} + \redoval{({\eta^{2}\omega^{4}\mu})} & b_{1} + \eta\omega^{1}\mu & c_{1} + \bluerect{({\eta\omega^{2}\mu})} \\
a_{2} + \omega^{2}\mu & b_{2} + \bluerect{({\eta^{2}\omega^{1}\mu})} & c_{2} + \eta^{2}\omega^{2}\mu \\
a_{3} + \omega^{3}\mu & b_{3} + \eta\omega^{3}\mu & c_{3} + \eta^{2}\omega^{3}\mu \\
a_{4} + \redoval{({\omega^{1}\mu})} & b_{4} + \eta\omega^{4}\mu & c_{4} + \redoval{({\omega^{4}\mu})}
\end{pmatrix}.
\end{array}
\end{equation}

The permutations needed are part of the problem statement, hence
globally known. So in this example, both parties are going to
interpolate cubic polynomials $\sigma_{a},\sigma_{b},\sigma_{c}$ that
encode the weird coefficients row-by-row:
\[
\begin{pmatrix}
\sigma_{a}( \omega^{1} ) = \redoval{({\eta^{2}\omega^{4}})} & \sigma_{b}( \omega^{1} ) = \eta\omega^{1} & \sigma_{c}( \omega^{1} ) = \bluerect{({\eta\omega^{2}})} \\
\sigma_{a}( \omega^{2} ) = \omega^{2} & \sigma_{b}( \omega^{2} ) = \bluerect{({\eta^{2}\omega^{1}})} & \sigma_{c}( \omega^{2} ) = \eta^{2}\omega^{2} \\
\sigma_{a}( \omega^{3} ) = \omega^{3} & \sigma_{b}( \omega^{3} ) = \eta\omega^{3} & \sigma_{c}( \omega^{3} ) = \eta^{2}\omega^{3} \\
\sigma_{a}( \omega^{4} ) = \redoval{(\omega^{1})} & \sigma_{b}( \omega^{4} ) = \eta\omega^{4} & \sigma_{c}( \omega^{4} ) = \redoval{(\omega^{4})}
\end{pmatrix}.
\]
Then the prover can start defining accumulator
polynomials, after re-introducing the random challenge $\lambda$ from
permutation check. We are going to need six in all, three for each side
of Equation \eqref{copy3}: We call them $F_{a}$, $F_{b}$,
$F_{c}$, $F'_{a}$, $F'_{b}$, $F'_{c}$. The ones on the left-hand
side are interpolated so that
\begin{equation}
\label{copycheck-left}
\begin{aligned}
F_{a}( \omega^{k} ) & = \prod_{i \leq k}( a_{i} + \omega^{i}\mu + \lambda ) \\
F_{b}( \omega^{k} ) & = \prod_{i \leq k}( b_{i} + \eta\omega^{i}\mu + \lambda ) \\
F_{c}( \omega^{k} ) & = \prod_{i \leq k}( c_{i} + \eta^{2}\omega^{i}\mu + \lambda ),
\end{aligned}
\end{equation}
while the ones on the right have the extra permutation polynomials
\begin{equation}
\label{copycheck-right}
\begin{aligned}
F'_{a}( \omega^{k} ) & = \prod_{i \leq k}( a_{i} + \sigma_{a}( \omega^{i} )\mu + \lambda ) \\
F'_{b}( \omega^{k} ) & = \prod_{i \leq k}( b_{i} + \sigma_{b}( \omega^{i} )\mu + \lambda ) \\
F'_{c}( \omega^{k} ) & = \prod_{i \leq k}( c_{i} + \sigma_{c}( \omega^{i} )\mu + \lambda ).
\end{aligned}
\end{equation}

And then we can run essentially the algorithm from before. There are six
initialization conditions
\begin{equation}
\label{copycheck-init}
\begin{aligned}
F_{a}( \omega^{1} ) & = A( \omega^{1} ) + \omega^{1}\mu + \lambda \\
F_{b}( \omega^{1} ) & = B( \omega^{1} ) + \eta\omega^{1}\mu + \lambda \\
F_{c}( \omega^{1} ) & = C( \omega^{1} ) + \eta^{2}\omega^{1}\mu + \lambda \\
F'_{a}( \omega^{1} ) & = A( \omega^{1} ) + \sigma_{a}( \omega^{1} )\mu + \lambda \\
F'_{b}( \omega^{1} ) & = B( \omega^{1} ) + \sigma_{b}( \omega^{1} )\mu + \lambda \\
F'_{c}( \omega^{1} ) & = C( \omega^{1} ) + \sigma_{c}( \omega^{1} )\mu + \lambda.
\end{aligned}
\end{equation}
and six accumulation conditions
\begin{equation}
\label{copycheck-accum}
\begin{aligned}
F_{a}(\omega X) & = F_{a}(X) \cdot ( A(\omega X) + X\mu + \lambda ) \\
F_{b}(\omega X) & = F_{b}(X) \cdot ( B(\omega X) + \eta X\mu + \lambda ) \\
F_{c}(\omega X) & = F_{c}(X) \cdot ( C(\omega X) + \eta^{2}X\mu + \lambda ) \\
F'_{a}(\omega X) & = F'_{a}(X) \cdot ( A(\omega X) + \sigma_{a}(X)\mu + \lambda ) \\
F'_{b}(\omega X) & = F'_{b}(X) \cdot ( B(\omega X) + \sigma_{b}(X)\mu + \lambda ) \\
F'_{c}(\omega X) & = F'_{c}(X) \cdot ( C(\omega X) + \sigma_{c}(X)\mu + \lambda )
\end{aligned}
\end{equation}
before the final product condition
\begin{equation}
\label{copycheck-final}
F_{a}(1)F_{b}(1)F_{c}(1) = F'_{a}(1)F'_{b}(1)F'_{c}(1).
\end{equation}

To summarize, the copy check goes as follows:

\begin{algo}{Copy check}{}
\leavevmode
\begin{enumerate}
\setcounter{enumi}{-1}
\item
  Peggy has already sent the three commitments
  \(\operatorname{Com}(A),\operatorname{Com}(B),\operatorname{Com}(C)\)
  to Victor; these commitments bind her to the values of all the
  variables $a_{i}$, $b_{i}$, and $c_{i}$.
\item
  Both parties compute the degree $n - 1$ polynomials
  $\sigma_{a},\sigma_{b},\sigma_{c} \in \FF_{q}[X]$
  described above, based on the copy constraints in the problem
  statement.
\item
  Victor chooses random challenges $\mu,\lambda \in \FF_{q}$
  and sends them to Peggy.
\item
  Peggy interpolates the six accumulator polynomials $F_{a}$, \ldots,
  $F'_{c}$ defined in Equation \eqref{copycheck-left}
  and Equation \eqref{copycheck-right}.
\item
  Peggy uses Algorithm \ref{tcb:root-check} to prove
  Equation \eqref{copycheck-init} holds.
\item
  Peggy uses Algorithm \ref{tcb:root-check} to prove
  Equation \eqref{copycheck-accum} holds for
  $X \in \{ \omega,\omega^{2},\ldots,\omega^{n - 1} \}$.
\item
  Peggy uses Algorithm \ref{tcb:root-check} to prove
  Equation \eqref{copycheck-final} holds.
\end{enumerate}
\end{algo}
